\chapter{Chapter09 — Learning Consequences Resources}
%%% generated by the blueprint pipeline from FABL.Chapter09.LearningConsequencesResources
\section{Overview}

Declarations of the Lean module \texttt{FABL.Chapter09.LearningConsequencesResources}.

\section{Declarations}

\begin{definition}[Subsquare-root learning regime]
\lean{FABL.LearningConsequencesSubsqrtRegime}
\label{FABL.LearningConsequencesSubsqrtRegime}
\leanok
\uses{FABL.PositiveInfluenceBound, FABL.PositiveLearningParameter}
Let $n,C,Q$ be natural numbers, let $t$ be a positive rational, and let $\varepsilon$ be
rational with $0<\varepsilon\le 1/2$. The parameters lie in the subsquare-root regime
when $1\le n$, $1/(Q+1)\le\varepsilon$, and $t\le
(C+1)(\sqrt{\lceil\log_2(n+1)\rceil}+1)$, where the logarithmic quantity is rounded up
to an integer before taking its square root.
\end{definition}

\begin{definition}[Learning degree coefficient]
\lean{FABL.learningConsequencesDegreeCoefficient}
\label{FABL.learningConsequencesDegreeCoefficient}
\leanok
For natural numbers $C$ and $Q$, define the learning degree coefficient
$\kappa(C,Q)=8(C+1)(Q+1)+1$.
\end{definition}

\begin{definition}[Learning consequences coordinate coefficient]
\lean{FABL.learningConsequencesCoordinateCoefficient}
\label{FABL.learningConsequencesCoordinateCoefficient}
\leanok
\uses{FABL.learningConsequencesDegreeCoefficient}
For natural numbers $C$ and $Q$, define the coordinate coefficient by $128(Q+1)^2(C+1)^3
+ 4K(C,Q) + 5$, where $K(C,Q)$ is the degree coefficient used in the learning
consequences bound.
\end{definition}

\begin{definition}[Family exponent]
\lean{FABL.learningConsequencesFamilyExponent}
\label{FABL.learningConsequencesFamilyExponent}
\leanok
\uses{FABL.learningConsequencesCoordinateCoefficient, FABL.learningConsequencesDegreeCoefficient}
For natural numbers $C$ and $Q$, define the family exponent to be twice the degree
coefficient times one more than the coordinate coefficient; explicitly, it is
$2\bigl(8(C+1)(Q+1)+1\bigr)\bigl(128(Q+1)^2(C+1)^3+4(8(C+1)(Q+1)+1)+6\bigr)$.
\end{definition}

\begin{theorem}[Successor bounded by a power of two]
\lean{FABL.nat_succ_le_two_pow}
\label{FABL.nat_succ_le_two_pow}
\leanok
\difficulty{4}
For every natural number $m$, one has $m+1\le2^m$.
\end{theorem}

\begin{theorem}[A natural number bounded by its power of two]
\lean{FABL.nat_self_le_two_pow}
\label{FABL.nat_self_le_two_pow}
\leanok
\uses{FABL.nat_succ_le_two_pow}
\difficulty{1}
For every natural number $m$, one has $m\le2^m$.
\end{theorem}

\begin{theorem}[Power bound for the binary ceiling logarithm]
\lean{FABL.two_pow_clog_le_two_mul}
\label{FABL.two_pow_clog_le_two_mul}
\leanok
\difficulty{4}
For every positive natural number $x$, the power $2^{\\operatorname{clog}_2(x)}$ is at
most $2x$, where $\operatorname{clog}_2(x)$ is the least natural number $k$ such that
$x\le 2^k$.
\end{theorem}

\begin{theorem}[Degree bound in the subsquare-root regime]
\lean{FABL.learningConsequencesDegree_le_subsqrt}
\label{FABL.learningConsequencesDegree_le_subsqrt}
\leanok
\uses{FABL.LearningConsequencesSubsqrtRegime, FABL.PositiveInfluenceBound, FABL.PositiveLearningParameter, FABL.learningConsequencesDegree, FABL.learningConsequencesDegreeCoefficient}
\difficulty{5}
Let $n,C,Q$ be nonnegative integers, let $t>0$, and let $0<\varepsilon\leq 1/2$. Suppose
$n\geq1$, $1/(Q+1)\leq\varepsilon$, and $t\leq(C+1)(\sqrt{\lceil\log_2(n+1)\rceil}+1)$.
For $d=\lceil 8t/\varepsilon\rceil$, one has
$d\leq[8(C+1)(Q+1)+1](\sqrt{\lceil\log_2(n+1)\rceil}+1)$.
\end{theorem}

\begin{theorem}[Natural Square Root Bound]
\lean{FABL.sqrt_add_one_sq_le_two_mul_add_one}
\label{FABL.sqrt_add_one_sq_le_two_mul_add_one}
\leanok
\difficulty{3}
For every natural number $L$, $(\lfloor\sqrt{L}\rfloor+1)^2\le 2(L+1)$, where
$\lfloor\sqrt{L}\rfloor$ denotes the integer square root of $L$.
\end{theorem}

\begin{theorem}[Coordinate ratio bound]
\lean{FABL.learningConsequencesCoordinateRatio_le_two_pow}
\label{FABL.learningConsequencesCoordinateRatio_le_two_pow}
\leanok
\uses{FABL.LearningConsequencesSubsqrtRegime, FABL.PositiveInfluenceBound, FABL.PositiveLearningParameter, FABL.learningConsequencesDegree, FABL.learningConsequencesDegreeCoefficient, FABL.learningConsequencesDegree_le_subsqrt, FABL.learningConsequencesInternalError, FABL.learningConsequencesInternalError_pos, FABL.learningConsequencesThreshold, FABL.nat_self_le_two_pow}
\difficulty{6}
Let $n,C,Q$ be nonnegative integers, let $t>0$, and let $0<\varepsilon\leq 1/2$. Suppose
$n\geq 1$, $1/(Q+1)\leq\varepsilon$, and $t\leq(C+1)(\sqrt{\lceil\log_2(n+1)\rceil}+1)$.
Define $d=\lceil 8t/\varepsilon\rceil$ and $\theta=(\varepsilon/8)^2t^{-2}(1/9)^d$. Then
$2t/\theta\leq
2^{(128(Q+1)^2(C+1)^3+4[8(C+1)(Q+1)+1]+3)(\sqrt{\lceil\log_2(n+1)\rceil}+1)}$.
\end{theorem}

\begin{theorem}[Coordinate cap bound]
\lean{FABL.learningConsequencesCoordinateCap_add_one_le_two_pow}
\label{FABL.learningConsequencesCoordinateCap_add_one_le_two_pow}
\leanok
\uses{FABL.LearningConsequencesSubsqrtRegime, FABL.PositiveInfluenceBound, FABL.PositiveLearningParameter, FABL.learningConsequencesCoordinateCap, FABL.learningConsequencesCoordinateCoefficient, FABL.learningConsequencesCoordinateRatio_le_two_pow, FABL.learningConsequencesDegreeCoefficient, FABL.learningConsequencesThreshold}
\difficulty{4}
Let $n\ge1$, let $C,Q$ be nonnegative integers, let $t>0$ be rational, and let
$0<\varepsilon\le1/2$ be rational. Suppose $1/(Q+1)\le\varepsilon$ and
$t\le(C+1)(\sqrt{\lceil\log_2(n+1)\rceil}+1)$. Define $d=\lceil 8t/\varepsilon\rceil$
and $T=(\varepsilon/8)^2t^{-2}(1/9)^d$. Then the coordinate cap $\lceil 2t/T\rceil$,
increased by one, is at most $2^{K(\sqrt{\lceil\log_2(n+1)\rceil}+1)}$, where
$K=128(Q+1)^2(C+1)^3+4(8(C+1)(Q+1)+1)+5$.
\end{theorem}

\begin{theorem}[Fixed-power bound for the family size]
\lean{FABL.learningConsequencesFamilyBound_le_fixed_power}
\label{FABL.learningConsequencesFamilyBound_le_fixed_power}
\leanok
\uses{FABL.LearningConsequencesSubsqrtRegime, FABL.PositiveInfluenceBound, FABL.PositiveLearningParameter, FABL.learningConsequencesCoordinateCap, FABL.learningConsequencesCoordinateCap_add_one_le_two_pow, FABL.learningConsequencesCoordinateCoefficient, FABL.learningConsequencesDegree, FABL.learningConsequencesDegreeCoefficient, FABL.learningConsequencesDegree_le_subsqrt, FABL.learningConsequencesFamilyBound, FABL.learningConsequencesFamilyExponent, FABL.nat_succ_le_two_pow, FABL.sqrt_add_one_sq_le_two_mul_add_one, FABL.two_pow_clog_le_two_mul}
\difficulty{5}
Let $n,C,Q$ be nonnegative integers, let $t>0$, and let $0<\varepsilon\leq1/2$. Suppose
$n\geq1$, $1/(Q+1)\leq\varepsilon$, and $t\leq(C+1)(\sqrt{\lceil\log_2(n+1)\rceil}+1)$.
Put $a=8(C+1)(Q+1)+1$, $b=128(Q+1)^2(C+1)^3+4a+5$, $e=2a(b+1)$,
$d=\lceil8t/\varepsilon\rceil$, $\theta=(\varepsilon/8)^2t^{-2}(1/9)^d$, and
$k=\lceil2t/\theta\rceil$. Then $(d+1)(k+1)^d\leq2^{2e}(n+1)^e$.
\end{theorem}

\begin{definition}[Learning-consequences family constant]
\lean{FABL.learningConsequencesFamilyConstant}
\label{FABL.learningConsequencesFamilyConstant}
\leanok
\uses{FABL.learningConsequencesFamilyExponent}
For natural numbers $C$ and $Q$, let $e(C,Q)$ denote the learning-consequences family
exponent. The family constant is $2^{2e(C,Q)}$.
\end{definition}

\begin{definition}[Learning consequences query exponent]
\lean{FABL.learningConsequencesQueryExponent}
\label{FABL.learningConsequencesQueryExponent}
\leanok
\uses{FABL.learningConsequencesFamilyExponent}
For natural numbers $C$ and $Q$, define the learning consequences query exponent to be
$8$ times the learning consequences family exponent at $C$ and $Q$, plus $2$.
\end{definition}

\begin{definition}[Learning consequences work exponent]
\lean{FABL.learningConsequencesWorkExponent}
\label{FABL.learningConsequencesWorkExponent}
\leanok
\uses{FABL.learningConsequencesFamilyExponent}
For natural numbers $C$ and $Q$, define the work exponent to be $8$ times the family
exponent for $(C,Q)$, plus $3$.
\end{definition}

\begin{definition}[Learning consequences query constant]
\lean{FABL.learningConsequencesQueryConstant}
\label{FABL.learningConsequencesQueryConstant}
\leanok
\uses{FABL.learningConsequencesFamilyConstant}
For natural numbers $C$ and $Q$, define the learning consequences query constant to be
$2^{40}\bigl(K(C,Q)+1\bigr)^8(Q+1)^{10}$, where $K(C,Q)$ denotes the learning
consequences family constant with parameters $C$ and $Q$.
\end{definition}

\begin{definition}[Learning consequences work constant]
\lean{FABL.learningConsequencesWorkConstant}
\label{FABL.learningConsequencesWorkConstant}
\leanok
\uses{FABL.learningConsequencesFamilyConstant}
For natural numbers $C$ and $Q$, define the learning consequences work constant to be
$2^{42}$ times $(K+1)^8(Q+1)^{10}$, where $K$ is the learning consequences family
constant for $C$ and $Q$.
\end{definition}

\begin{theorem}[Family bound with an additive one]
\lean{FABL.learningConsequencesFamilyBound_add_one_le_fixed_power}
\label{FABL.learningConsequencesFamilyBound_add_one_le_fixed_power}
\leanok
\uses{FABL.LearningConsequencesSubsqrtRegime, FABL.PositiveInfluenceBound, FABL.PositiveLearningParameter, FABL.learningConsequencesFamilyBound, FABL.learningConsequencesFamilyBound_le_fixed_power, FABL.learningConsequencesFamilyConstant, FABL.learningConsequencesFamilyExponent}
\difficulty{4}
Let $n,C,Q$ be nonnegative integers, let $t>0$, and let $0<\varepsilon\leq1/2$. Suppose
$n\geq1$, $1/(Q+1)\leq\varepsilon$, and $t\leq(C+1)(\sqrt{\lceil\log_2(n+1)\rceil}+1)$.
Put $a=8(C+1)(Q+1)+1$, $b=128(Q+1)^2(C+1)^3+4a+5$, $e=2a(b+1)$,
$d=\lceil8t/\varepsilon\rceil$, $\theta=(\varepsilon/8)^2t^{-2}(1/9)^d$, and
$k=\lceil2t/\theta\rceil$. Then $(d+1)(k+1)^d+1\leq(2^{2e}+1)(n+1)^e$.
\end{theorem}

\begin{theorem}[Power bound from an accuracy lower bound]
\lean{FABL.one_div_accuracy_pow_le}
\label{FABL.one_div_accuracy_pow_le}
\leanok
\uses{FABL.LearningConsequencesSubsqrtRegime, FABL.PositiveInfluenceBound, FABL.PositiveLearningParameter}
\difficulty{4}
Let $n,C,Q,r$ be natural numbers, let $t$ be a positive rational, and let $\varepsilon$
be a rational with $0<\varepsilon\le 1/2$. Suppose that $1\le n$, that
$1/(Q+1)\le\varepsilon$, and that $t\le (C+1)(\sqrt{\lceil\log_2(n+1)\rceil}+1)$. Then
$(1/\varepsilon)^r\le (Q+1)^r$.
\end{theorem}

\begin{theorem}[Resource expression power bound]
\lean{FABL.kmResourceExpression_le_fixed_power}
\label{FABL.kmResourceExpression_le_fixed_power}
\leanok
\difficulty{4}
Let $n$, $L$, $M$, $F$, $E$, and $Q$ be natural numbers, and let $\varepsilon$ be a
positive rational. If $M+1\le (F+1)(n+1)^E$ and $(1/\varepsilon)^{10}\le Q^{10}$, then
$L(n+1)^d(M+1)^8/\varepsilon^{10}\le L(F+1)^8Q^{10}(n+1)^{8E+d}$, where $d$ denotes the
natural number named `dimensionPower` and $L$ the natural number named `leading`.
\end{theorem}

\begin{theorem}[Polynomial query-budget bound]
\lean{FABL.boundedInfluenceQueryBudget_le_fixed_power}
\label{FABL.boundedInfluenceQueryBudget_le_fixed_power}
\leanok
\uses{FABL.LearningConsequencesSubsqrtRegime, FABL.PositiveInfluenceBound, FABL.PositiveLearningParameter, FABL.kmResourceExpression_le_fixed_power, FABL.learningConsequencesFamilyBound, FABL.learningConsequencesFamilyBound_add_one_le_fixed_power, FABL.learningConsequencesFamilyConstant, FABL.learningConsequencesFamilyExponent, FABL.learningConsequencesQueryConstant, FABL.learningConsequencesQueryExponent, FABL.one_div_accuracy_pow_le}
\difficulty{1}
Let $n,C,Q$ be natural numbers, let $t>0$ be rational, and let $0<\varepsilon\le1/2$ be
rational. Assume $1\le n$, $1/(Q+1)\le\varepsilon$, and
$t\le(C+1)(\sqrt{\lceil\log_2(n+1)\rceil}+1)$. Define $d=\lceil8t/\varepsilon\rceil$,
$\theta=(\varepsilon/8)^2t^{-2}(1/9)^d$, $c=\lceil2t/\theta\rceil$, $B=(d+1)(c+1)^d$,
$a=8(C+1)(Q+1)+1$, $b=128(Q+1)^2(C+1)^3+4a+5$, and $e=2a(b+1)$. Then
$2^{40}(n+1)^2(B+1)^8/\varepsilon^{10}\le 2^{40}(2^{2e}+1)^8(Q+1)^{10}(n+1)^{8e+2}$.
\end{theorem}

\begin{theorem}[Polynomial work bound in the sub-square-root regime]
\lean{FABL.boundedInfluenceWorkBudget_le_fixed_power}
\label{FABL.boundedInfluenceWorkBudget_le_fixed_power}
\leanok
\uses{FABL.LearningConsequencesSubsqrtRegime, FABL.PositiveInfluenceBound, FABL.PositiveLearningParameter, FABL.kmResourceExpression_le_fixed_power, FABL.learningConsequencesFamilyBound, FABL.learningConsequencesFamilyBound_add_one_le_fixed_power, FABL.learningConsequencesFamilyConstant, FABL.learningConsequencesFamilyExponent, FABL.learningConsequencesWorkConstant, FABL.learningConsequencesWorkExponent, FABL.one_div_accuracy_pow_le}
\difficulty{4}
Let $n,C,Q$ be natural numbers, let $t>0$ be rational, and let $0<\epsilon\le 1/2$ be
rational. Suppose $n\ge1$, $1/(Q+1)\le\epsilon$, and
$t\le(C+1)(\sqrt{\operatorname{clog}_2(n+1)}+1)$. Then
$2^{42}(n+1)^3(B+1)^8/\epsilon^{10}\le K(n+1)^E$, where $B$ is the family bound
determined by $t$ and $\epsilon$, $K$ is the work constant determined by $C,Q$, and $E$
is the work exponent determined by $C,Q$.
\end{theorem}

\begin{theorem}[Failure-bit count for bounded family confidence]
\lean{FABL.boundedFamilyCoefficientFailureBits_eq}
\label{FABL.boundedFamilyCoefficientFailureBits_eq}
\leanok
\uses{FABL.boundedFamilyCoefficientConfidence, FABL.fourierEstimatorFailureBits}
\difficulty{3}
For every positive integer $M$, if the confidence parameter is $1/(20M)$ and the
failure-bit count is the base-two ceiling logarithm of the ceiling of twice the
reciprocal of that parameter, then the count equals $\lceil\log_2(40M)\rceil$, with the
ceiling logarithm understood as the least integer $k$ such that $2^k$ is at least its
argument.
\end{theorem}

\begin{theorem}[Bound on bounded-family sample count]
\lean{FABL.boundedFamilySampleCount_cast_le}
\label{FABL.boundedFamilySampleCount_cast_le}
\leanok
\uses{FABL.PositiveLearningParameter, FABL.boundedFamilyCoefficientAccuracy, FABL.boundedFamilyCoefficientConfidence, FABL.boundedFamilyCoefficientFailureBits_eq, FABL.boundedFamilySampleCount, FABL.fourierEstimatorFailureBits, FABL.fourierEstimatorSampleCount_cast_le, FABL.nat_self_le_two_pow}
\difficulty{4}
Let $0<\varepsilon\le 1/2$ be rational and let $M$ be a positive integer. The
bounded-family sample count, formed using accuracy $\varepsilon/(2M)$ and confidence
$1/(20M)$, is at most $640M^3/\varepsilon^2$.
\end{theorem}

\begin{definition}[Second-stage sample constant]
\lean{FABL.learningConsequencesSecondStageSampleConstant}
\label{FABL.learningConsequencesSecondStageSampleConstant}
\leanok
\uses{FABL.learningConsequencesFamilyConstant}
For natural numbers $C$ and $Q$, the second-stage sample constant is $640$ times the
cube of the family constant at $(C,Q)$ times $(Q+1)^2$.
\end{definition}

\begin{definition}[Second-stage sample exponent]
\lean{FABL.learningConsequencesSecondStageSampleExponent}
\label{FABL.learningConsequencesSecondStageSampleExponent}
\leanok
\uses{FABL.learningConsequencesFamilyExponent}
For natural numbers $C$ and $Q$, the second-stage sample exponent is three times the
family exponent at $(C,Q)$.
\end{definition}

\begin{theorem}[Cubic Sample Expression Bound]
\lean{FABL.cubicSampleExpression_le_fixed_power}
\label{FABL.cubicSampleExpression_le_fixed_power}
\leanok
\difficulty{4}
Let $n,M,A,E,Q$ be natural numbers and let $\varepsilon$ be a positive rational. If
$M\le A(n+1)^E$ and $(1/\varepsilon)^2\le Q^2$, then $640M^3/\varepsilon^2\le
640A^3Q^2(n+1)^{3E}$.
\end{theorem}

\begin{theorem}[Polynomial sample bound in the dimension]
\lean{FABL.boundedFamilySampleCount_le_fixed_power}
\label{FABL.boundedFamilySampleCount_le_fixed_power}
\leanok
\uses{FABL.LearningConsequencesSubsqrtRegime, FABL.PositiveInfluenceBound, FABL.PositiveLearningParameter, FABL.boundedFamilySampleCount, FABL.boundedFamilySampleCount_cast_le, FABL.cubicSampleExpression_le_fixed_power, FABL.learningConsequencesFamilyBound, FABL.learningConsequencesFamilyBound_le_fixed_power, FABL.learningConsequencesFamilyBound_pos, FABL.learningConsequencesFamilyConstant, FABL.learningConsequencesFamilyExponent, FABL.learningConsequencesSecondStageSampleConstant, FABL.learningConsequencesSecondStageSampleExponent, FABL.one_div_accuracy_pow_le}
\difficulty{4}
Let $n,C,Q$ be natural numbers, let $t>0$ be rational, and let $0<\varepsilon\le 1/2$ be
rational. Assume $1\le n$, $1/(Q+1)\le\varepsilon$, and
$t\le(C+1)(\sqrt{\lceil\log_2(n+1)\rceil}+1)$. Form the family size bound
$B=(d+1)(c+1)^d$, where $d=\lceil 8t/\varepsilon\rceil$,
$\theta=(\varepsilon/8)^2t^{-2}(1/9)^d$, and $c=\lceil2t/\theta\rceil$. Then the
bounded-family sample count with family bound $B$ is at most $640\,2^{6e} (Q+1)^2
(n+1)^{3e}$, where $e=2a(b+1)$, $a=8(C+1)(Q+1)+1$, and $b=128(Q+1)^2(C+1)^3+4a+5$.
\end{theorem}

\begin{definition}[Singleton exponent coefficient]
\lean{FABL.learningConsequencesSingletonExponentCoefficient}
\label{FABL.learningConsequencesSingletonExponentCoefficient}
\leanok
\uses{FABL.learningConsequencesDegreeCoefficient}
For natural numbers $C$ and $Q$, the singleton exponent coefficient is $(8(C+1)(Q+1))^4$
plus eight times the degree coefficient at $(C,Q)$ plus four.
\end{definition}

\begin{theorem}[Rational bound on the inverse learning threshold]
\lean{FABL.one_div_learningConsequencesThreshold_sq_le_expression}
\label{FABL.one_div_learningConsequencesThreshold_sq_le_expression}
\leanok
\uses{FABL.LearningConsequencesSubsqrtRegime, FABL.PositiveInfluenceBound, FABL.PositiveLearningParameter, FABL.learningConsequencesDegree, FABL.learningConsequencesDegreeCoefficient, FABL.learningConsequencesDegree_le_subsqrt, FABL.learningConsequencesInternalError, FABL.learningConsequencesInternalError_pos, FABL.learningConsequencesThreshold}
\difficulty{6}
Let $n,C,Q$ be natural numbers, let $t$ be a positive rational, and let $\varepsilon$ be
rational with $0<\varepsilon\le 1/2$. Assume $1\le n$, $1/(Q+1)\le\varepsilon$, and
$t\le(C+1)(\sqrt{\lceil\log_2(n+1)\rceil}+1)$. Define $d=\lceil 8t/\varepsilon\rceil$
and $\tau=(\varepsilon/8)^2t^{-2}(1/9)^d$. Then $\tau^{-2}\le [
(8(C+1)(Q+1))^4(\sqrt{\lceil\log_2(n+1)\rceil}+1)^4
9^{2(8(C+1)(Q+1)+1)(\sqrt{\lceil\log_2(n+1)\rceil}+1)} ]$, with the right-hand side
interpreted as a rational number.
\end{theorem}

\begin{theorem}[Exponential bound for a singleton expression]
\lean{FABL.learningConsequencesSingletonExpression_le_two_pow}
\label{FABL.learningConsequencesSingletonExpression_le_two_pow}
\leanok
\uses{FABL.learningConsequencesDegreeCoefficient, FABL.learningConsequencesSingletonExponentCoefficient, FABL.nat_self_le_two_pow}
\difficulty{5}
For all nonnegative integers $n,C,Q$, let $a=8(C+1)(Q+1)+1$ and
$b=(8(C+1)(Q+1))^4+8a+4$. Then $(8(C+1)(Q+1))^4(\sqrt{\lceil\log_2(n+1)\rceil}+1)^4
9^{2a(\sqrt{\lceil\log_2(n+1)\rceil}+1)}\leq2^{b(\sqrt{\lceil\log_2(n+1)\rceil}+1)}$.
\end{theorem}

\begin{theorem}[Exponential bound on the inverse learning threshold]
\lean{FABL.one_div_learningConsequencesThreshold_sq_le_two_pow}
\label{FABL.one_div_learningConsequencesThreshold_sq_le_two_pow}
\leanok
\uses{FABL.LearningConsequencesSubsqrtRegime, FABL.PositiveInfluenceBound, FABL.PositiveLearningParameter, FABL.learningConsequencesSingletonExponentCoefficient, FABL.learningConsequencesSingletonExpression_le_two_pow, FABL.learningConsequencesThreshold, FABL.one_div_learningConsequencesThreshold_sq_le_expression}
\difficulty{1}
Let $n,C,Q$ be natural numbers, let $t$ be a positive rational, and let $\varepsilon$ be
rational with $0<\varepsilon\le 1/2$. Assume $1\le n$, $1/(Q+1)\le\varepsilon$, and
$t\le(C+1)(\sqrt{\lceil\log_2(n+1)\rceil}+1)$. Define $d=\lceil 8t/\varepsilon\rceil$
and $\tau=(\varepsilon/8)^2t^{-2}(1/9)^d$. Then $\tau^{-2}\le
2^{(((8(C+1)(Q+1))^4+8(8(C+1)(Q+1)+1)+4)(\sqrt{\lceil\log_2(n+1)\rceil}+1))}$, with the
right-hand side interpreted as a rational number.
\end{theorem}

\begin{theorem}[Singleton accuracy reciprocal bound]
\lean{FABL.one_div_singletonInfluenceAccuracy_sq_le}
\label{FABL.one_div_singletonInfluenceAccuracy_sq_le}
\leanok
\uses{FABL.LearningConsequencesSubsqrtRegime, FABL.PositiveInfluenceBound, FABL.PositiveLearningParameter, FABL.learningConsequencesSingletonExponentCoefficient, FABL.learningConsequencesThreshold, FABL.one_div_learningConsequencesThreshold_sq_le_two_pow, FABL.singletonInfluenceAccuracy, FABL.singletonInfluenceAccuracy_value}
\difficulty{4}
Let $n\geq 0$ and $C,Q\in\mathbb{N}$. Let $t$ be a positive rational number and let
$\varepsilon$ be a rational number with $0<\varepsilon\leq 1/2$. Assume $1\leq n$,
$1/(Q+1)\leq\varepsilon$, and $t\leq(C+1)(\sqrt{\lceil\log_2(n+1)\rceil}+1)$. Define
$d=\lceil 8t/\varepsilon\rceil$, $T=(\varepsilon/8)^2/t^2\,(1/9)^d$, and
$a=\min(T/4,1/2)$. Then $a^{-2}\leq 16\cdot
2^{(((8(C+1)(Q+1))^4+8(8(C+1)(Q+1)+1)+4)(\sqrt{\lceil\log_2(n+1)\rceil}+1))}$.
\end{theorem}

\begin{theorem}[Failure bits for the confidence parameter]
\lean{FABL.singletonInfluenceFailureBits_eq}
\label{FABL.singletonInfluenceFailureBits_eq}
\leanok
\uses{FABL.fourierEstimatorFailureBits, FABL.singletonInfluenceConfidence}
\difficulty{3}
For every natural number $n$, if the confidence parameter is $1/(20(n+1))$ and the
number of failure bits at parameter $\delta$ is $\operatorname{clog}_2\lceil
2/\delta\rceil$, then that number is $\operatorname{clog}_2(40(n+1))$.
\end{theorem}

\begin{theorem}[Subsquare-root sample-count bound]
\lean{FABL.singletonInfluenceSampleCount_cast_le_subsqrt}
\label{FABL.singletonInfluenceSampleCount_cast_le_subsqrt}
\leanok
\uses{FABL.LearningConsequencesSubsqrtRegime, FABL.PositiveInfluenceBound, FABL.PositiveLearningParameter, FABL.fourierEstimatorFailureBits, FABL.fourierEstimatorSampleCount_cast_le, FABL.learningConsequencesSingletonExponentCoefficient, FABL.nat_self_le_two_pow, FABL.one_div_singletonInfluenceAccuracy_sq_le, FABL.singletonInfluenceAccuracy, FABL.singletonInfluenceConfidence, FABL.singletonInfluenceFailureBits_eq, FABL.singletonInfluenceSampleCount}
\difficulty{5}
Let $n,C,Q$ be natural numbers, let $t$ be a positive rational, and let $\varepsilon$ be
rational with $0<\varepsilon\le 1/2$. Suppose $1\le n$, $1/(Q+1)\le\varepsilon$, and
$t\le(C+1)(\sqrt{\operatorname{clog}_2(n+1)}+1)$. Set
$E=(8(C+1)(Q+1))^4+8(8(C+1)(Q+1)+1)+4$. Then the sample count, viewed as a rational
number, is at most $2560(n+1)2^{E(\sqrt{\operatorname{clog}_2(n+1)}+1)}$.
\end{theorem}

\begin{definition}[Singleton sample constant]
\lean{FABL.learningConsequencesSingletonSampleConstant}
\label{FABL.learningConsequencesSingletonSampleConstant}
\leanok
\uses{FABL.learningConsequencesSingletonExponentCoefficient}
For natural numbers $C$ and $Q$, the singleton sample constant is $2560$ times $2$
raised to twice the singleton exponent coefficient at $(C,Q)$.
\end{definition}

\begin{definition}[Singleton sample exponent]
\lean{FABL.learningConsequencesSingletonSampleExponent}
\label{FABL.learningConsequencesSingletonSampleExponent}
\leanok
\uses{FABL.learningConsequencesSingletonExponentCoefficient}
For natural numbers $C$ and $Q$, the singleton sample exponent is the singleton exponent
coefficient at $(C,Q)$ plus one.
\end{definition}

\begin{theorem}[A bound on a logarithmic exponent]
\lean{FABL.two_pow_mul_sqrt_clog_add_one_le_fixed_power}
\label{FABL.two_pow_mul_sqrt_clog_add_one_le_fixed_power}
\leanok
\uses{FABL.two_pow_clog_le_two_mul}
\difficulty{4}
For all nonnegative integers $n$ and $H$, let $r$ be the least nonnegative integer such
that $n+1\le 2^r$. Then $2^{H(\lfloor\sqrt r\rfloor+1)}\le 2^{2H}(n+1)^H$.
\end{theorem}

\begin{theorem}[Polynomial sample-count bound]
\lean{FABL.singletonInfluenceSampleCount_le_fixed_power}
\label{FABL.singletonInfluenceSampleCount_le_fixed_power}
\leanok
\uses{FABL.LearningConsequencesSubsqrtRegime, FABL.PositiveInfluenceBound, FABL.PositiveLearningParameter, FABL.learningConsequencesSingletonExponentCoefficient, FABL.learningConsequencesSingletonSampleConstant, FABL.learningConsequencesSingletonSampleExponent, FABL.singletonInfluenceSampleCount, FABL.singletonInfluenceSampleCount_cast_le_subsqrt, FABL.two_pow_mul_sqrt_clog_add_one_le_fixed_power}
\difficulty{4}
Let $n,C,Q$ be natural numbers, let $t$ be a positive rational, and let $\varepsilon$ be
rational with $0<\varepsilon\le 1/2$. Suppose $1\le n$, $1/(Q+1)\le\varepsilon$, and
$t\le(C+1)(\sqrt{\operatorname{clog}_2(n+1)}+1)$. Set
$E=(8(C+1)(Q+1))^4+8(8(C+1)(Q+1)+1)+4$. Then the sample count is at most $2560\cdot
2^{2E}(n+1)^{E+1}$.
\end{theorem}

\begin{definition}[Random-example exponent]
\lean{FABL.learningConsequencesRandomExampleExponent}
\label{FABL.learningConsequencesRandomExampleExponent}
\leanok
\uses{FABL.learningConsequencesFamilyExponent, FABL.learningConsequencesSecondStageSampleExponent, FABL.learningConsequencesSingletonSampleExponent}
For natural numbers $C$ and $Q$, the random-example exponent is the maximum of the
singleton sample exponent at $(C,Q)$ plus one and the sum of the family exponent and
second-stage sample exponent at $(C,Q)$.
\end{definition}

\begin{definition}[Random-example constant]
\lean{FABL.learningConsequencesRandomExampleConstant}
\label{FABL.learningConsequencesRandomExampleConstant}
\leanok
\uses{FABL.learningConsequencesFamilyConstant, FABL.learningConsequencesSecondStageSampleConstant, FABL.learningConsequencesSingletonSampleConstant}
For natural numbers $C$ and $Q$, the random-example constant is the singleton sample
constant at $(C,Q)$ plus the product of the family constant and second-stage sample
constant at $(C,Q)$.
\end{definition}

\begin{definition}[Random-example work exponent]
\lean{FABL.learningConsequencesRandomExampleWorkExponent}
\label{FABL.learningConsequencesRandomExampleWorkExponent}
\leanok
\uses{FABL.learningConsequencesFamilyExponent, FABL.learningConsequencesSecondStageSampleExponent, FABL.learningConsequencesSingletonSampleExponent}
For natural numbers $C$ and $Q$, the random-example work exponent is the maximum of the
singleton sample exponent at $(C,Q)$ plus two and the sum of the family exponent,
second-stage sample exponent, and one at $(C,Q)$.
\end{definition}

\begin{definition}[Random example work constant]
\lean{FABL.learningConsequencesRandomExampleWorkConstant}
\label{FABL.learningConsequencesRandomExampleWorkConstant}
\leanok
\uses{FABL.learningConsequencesFamilyConstant, FABL.learningConsequencesSecondStageSampleConstant, FABL.learningConsequencesSingletonSampleConstant}
For natural numbers $C$ and $Q$, define the random-example work constant to be $2$ times
the singleton-sample constant for $C,Q$, plus $1$, plus $2$ times the family constant
for $C,Q$ multiplied by the second-stage sample constant for $C,Q$.
\end{definition}

\begin{theorem}[Polynomial bound on example budget]
\lean{FABL.monotoneBoundedInfluenceRandomExampleBudget_le_fixed_power}
\label{FABL.monotoneBoundedInfluenceRandomExampleBudget_le_fixed_power}
\leanok
\uses{FABL.LearningConsequencesSubsqrtRegime, FABL.PositiveInfluenceBound, FABL.PositiveLearningParameter, FABL.boundedFamilySampleCount, FABL.boundedFamilySampleCount_le_fixed_power, FABL.learningConsequencesFamilyBound, FABL.learningConsequencesFamilyBound_le_fixed_power, FABL.learningConsequencesFamilyBound_pos, FABL.learningConsequencesFamilyConstant, FABL.learningConsequencesFamilyExponent, FABL.learningConsequencesRandomExampleConstant, FABL.learningConsequencesRandomExampleExponent, FABL.learningConsequencesSecondStageSampleConstant, FABL.learningConsequencesSecondStageSampleExponent, FABL.learningConsequencesSingletonSampleConstant, FABL.learningConsequencesSingletonSampleExponent, FABL.monotoneBoundedInfluenceRandomExampleBudget, FABL.singletonInfluenceSampleCount, FABL.singletonInfluenceSampleCount_le_fixed_power}
\difficulty{5}
Let $n\ge1$, let $t>0$ be rational, let $0<\varepsilon\le\tfrac12$ be rational, and let
$C,Q$ be nonnegative integers. Suppose $\varepsilon\ge 1/(Q+1)$ and
$t\le(C+1)\bigl(\sqrt{\lceil\log_2(n+1)\rceil}+1\bigr)$. Then the monotone
bounded-influence random-example budget is at most $K(C,Q)(n+1)^{E(C,Q)}$, where
$K(C,Q)$ and $E(C,Q)$ are the specified random-example constant and exponent.
\end{theorem}

\begin{theorem}[Polynomial bound on work budget]
\lean{FABL.monotoneBoundedInfluenceWorkBudget_le_fixed_power}
\label{FABL.monotoneBoundedInfluenceWorkBudget_le_fixed_power}
\leanok
\uses{FABL.LearningConsequencesSubsqrtRegime, FABL.PositiveInfluenceBound, FABL.PositiveLearningParameter, FABL.boundedFamilySampleCount, FABL.boundedFamilySampleCount_le_fixed_power, FABL.learningConsequencesFamilyBound, FABL.learningConsequencesFamilyBound_le_fixed_power, FABL.learningConsequencesFamilyBound_pos, FABL.learningConsequencesFamilyConstant, FABL.learningConsequencesFamilyExponent, FABL.learningConsequencesRandomExampleWorkConstant, FABL.learningConsequencesRandomExampleWorkExponent, FABL.learningConsequencesSecondStageSampleConstant, FABL.learningConsequencesSecondStageSampleExponent, FABL.learningConsequencesSingletonSampleConstant, FABL.learningConsequencesSingletonSampleExponent, FABL.monotoneBoundedInfluenceWorkBudget, FABL.singletonInfluenceSampleCount, FABL.singletonInfluenceSampleCount_le_fixed_power}
\difficulty{5}
Let $n\ge1$, let $t>0$ and $0<\varepsilon\le\tfrac12$ be rational, and let $C,Q$ be
nonnegative integers. Suppose $\varepsilon\ge1/(Q+1)$ and
$t\le(C+1)\bigl(\sqrt{\lceil\log_2(n+1)\rceil}+1\bigr)$. Then the monotone
bounded-influence work budget is at most $W(C,Q)(n+1)^{F(C,Q)}$, where $W(C,Q)$ and
$F(C,Q)$ are the specified work constant and exponent.
\end{theorem}

\begin{theorem}[Polynomial bounds on program costs]
\lean{FABL.monotoneBoundedInfluenceRandomExampleProgram_cost_le_fixed_power}
\label{FABL.monotoneBoundedInfluenceRandomExampleProgram_cost_le_fixed_power}
\leanok
\uses{FABL.BooleanFunction, FABL.LearningConsequencesSubsqrtRegime, FABL.LearningCost, FABL.LearningProgram.runWithCost, FABL.PositiveInfluenceBound, FABL.PositiveLearningParameter, FABL.SparseFourierHypothesis, FABL.learningConsequencesRandomExampleConstant, FABL.learningConsequencesRandomExampleExponent, FABL.learningConsequencesRandomExampleWorkConstant, FABL.learningConsequencesRandomExampleWorkExponent, FABL.monotoneBoundedInfluenceRandomExampleBudget_le_fixed_power, FABL.monotoneBoundedInfluenceRandomExampleProgram, FABL.monotoneBoundedInfluenceRandomExampleProgram_cost_le, FABL.monotoneBoundedInfluenceWorkBudget_le_fixed_power}
\difficulty{3}
Let $n\ge1$, let $C,Q$ be nonnegative integers, let $t>0$ and $0<\varepsilon\le\tfrac12$
be rational, and suppose $\varepsilon\ge1/(Q+1)$ and
$t\le(C+1)\bigl(\sqrt{\lceil\log_2(n+1)\rceil}+1\bigr)$. For any Boolean target on the
$n$-dimensional sign cube, every outcome in the support of the program's cost
distribution uses at most $K(C,Q)(n+1)^{E(C,Q)}$ random examples, makes zero membership
queries, and incurs work at most $W(C,Q)(n+1)^{F(C,Q)}$, where $K,E,W,F$ are the
specified random-example constants and exponents.
\end{theorem}

\begin{theorem}[Decision tree influence bound in the sub-square-root regime]
\lean{FABL.decisionTreeSizeInfluenceBound_subsqrtRegime}
\label{FABL.decisionTreeSizeInfluenceBound_subsqrtRegime}
\leanok
\uses{FABL.LearningConsequencesSubsqrtRegime, FABL.PositiveLearningParameter, FABL.decisionTreeSizeInfluenceBound}
\difficulty{4}
Let $n,d,Q$ be natural numbers, let $$ be a rational learning parameter with $0<\langle
1/2$, and suppose $1\le n$ and $1/(Q+1)\le$. Then the decision-tree influence bound at
dimension $n$ and size parameter $d$ satisfies the sub-square-root regime conditions:
$1\le n$, $1/(Q+1)\le$, and
$\sqrt{d\,\lceil\log_2(n+1)\rceil}+1\le(d+2)(\sqrt{\lceil\log_2(n+1)\rceil}+1)$.
\end{theorem}

\begin{theorem}[Polynomial bound for the learning family]
\lean{FABL.decisionTreeSizeFamilyBound_le_fixed_power}
\label{FABL.decisionTreeSizeFamilyBound_le_fixed_power}
\leanok
\uses{FABL.PositiveLearningParameter, FABL.decisionTreeSizeInfluenceBound, FABL.decisionTreeSizeInfluenceBound_subsqrtRegime, FABL.learningConsequencesFamilyBound, FABL.learningConsequencesFamilyBound_le_fixed_power, FABL.learningConsequencesFamilyConstant, FABL.learningConsequencesFamilyExponent}
\difficulty{1}
Let $n,d,Q\in\mathbb{N}$ with $1\leq n$, and let $0<\varepsilon\leq 1/2$ satisfy
$1/(Q+1)\leq\varepsilon$. For the quantities defined by
$t=\sqrt{d\,\lceil\log_2(n+1)\rceil}+1$, $e=\lceil t/(\varepsilon/8)\rceil$,
$\theta=(\varepsilon/8)^2/t^2\,(1/9)^e$, $c=\lceil 2t/\theta\rceil$, $B=(e+1)(c+1)^e$,
$r=2e_0(c_0+1)$ with $e_0=8(d+2)(Q+1)+1$ and $c_0=128(Q+1)^2(d+2)^3+4e_0+5$, and
$C=2^r$, one has $B\leq C(n+1)^r$.
\end{theorem}

\begin{theorem}[Random-example and work bounds]
\lean{FABL.monotoneDecisionTreeSizeRandomExampleProgram_cost_le_fixed_power}
\label{FABL.monotoneDecisionTreeSizeRandomExampleProgram_cost_le_fixed_power}
\leanok
\uses{FABL.BooleanFunction, FABL.LearningCost, FABL.LearningProgram.runWithCost, FABL.PositiveLearningParameter, FABL.SparseFourierHypothesis, FABL.decisionTreeSizeInfluenceBound_subsqrtRegime, FABL.learningConsequencesRandomExampleConstant, FABL.learningConsequencesRandomExampleExponent, FABL.learningConsequencesRandomExampleWorkConstant, FABL.learningConsequencesRandomExampleWorkExponent, FABL.monotoneBoundedInfluenceRandomExampleProgram_cost_le_fixed_power, FABL.monotoneDecisionTreeSizeRandomExampleProgram}
\difficulty{3}
Let $n\ge 1$, let $d,Q$ be nonnegative integers, and let $0<\varepsilon\le 1/2$. For any
Boolean function on the sign cube and any outcome in the support of the program that
uses the decision-tree-size influence bound with parameters $n,d,\varepsilon$, the
number of random examples is at most $C(d+1,Q)(n+1)^{E(d+1,Q)}$, the number of
membership queries is zero, and the charged work is at most $W(d+1,Q)(n+1)^{F(d+1,Q)}$,
where $C,E,W,F$ are the specified learning-consequence constants and exponents for
random examples and work, respectively. The bound assumes $1/(Q+1)\le\varepsilon$.
\end{theorem}
